\chapter{Evidential Relevance}
\label{sec:p15-rationality-4-3}

Evidence can increase a hypothesis's probability without making it probable,
and a highly probable hypothesis can receive no support from a particular
observation. Arthur Merin's chapter~4.3 of \emph{The Handbook of Rationality}
develops this distinction through signed qualitative relevance, quantitative
measures, task-relative comparisons, and rules for compound evidence.
This treatment follows the 2021 chapter and its references
\parencite[pp.~265--275]{Merin2021Relevance}. Merin's argument for the
status of the axioms, the provenance of the measures, notes~1, 5, 8, 11--14,
17--20, the routes to the old-evidence problem, Pearl's diagnostic thesis,
measure pluralism, the comparison with Carnap's relevance theory, the
characterization of the measure Z, and the standing of every cited work are given in
Chapter~\ref{sec:p15-w02b-p15-024-relevance}.
Its models and examples are mathematical or philosophical unless expressly
identified as reports of psychological research. The probability model,
the update rule, and the evidential interpretation require separate warrant.

\section{Speaking for, speaking against, and neither}
\label{sec:p15-r43-qualitative}

Merin begins with dark clouds, rain, dry weather, and a rise in the Dow.
In a typical stipulated belief context, dark clouds favor rain, count
against dryness, and leave the market prediction unchanged. Learning
that there are no dark clouds would reverse the first two directions.
These are contextual judgments, not weather data or a measured
cloud--market independence. A source can speak for a target without
making it certain, and speak against it without making it impossible.

\begin{hbdefinition}{def:p15-r43-polarity}{Qualitative probabilistic relevance}
Work on an event algebra with a finitely additive probability \(P\).
Propositions are represented by their truth events; juxtaposition
\(AB\) denotes conjunction, and \(\neg A\) denotes complement.
Define the convenient covariance expression
\[
             \Delta_P(A,B)=P(AB)-P(A)P(B).
\]
The source relations \(A\,\mathrm{pos}_P\,B\), \((A\perp B)_P\), and
\(A\,\mathrm{neg}_P\,B\) mean respectively that this expression is
positive, zero, or negative. When \(P(A)>0\), these are exactly the
comparisons of \(P(B\mid A)\) with \(P(B)\). If \(P(A)=0\), the
ratio conditional is undefined but the relevance relation is neutral.

An update \(P_A=P(\cdot\mid A)\) represents learning just \(A\) with
certainty, under the conditioning premise. The conditional probability
\(P(B\mid A)\), its hypothetical update interpretation, and actual
observed learning are different objects.
\end{hbdefinition}

\begin{hbdefinition}{def:p15-r43-extreme}{Conclusive relevance and conditional irrelevance}
\(A\,\mathrm{expos}_P\,B\) means
\(P(A)>0\), \(P(B)<1\), and \(P(B\mid A)=1\).
\(A\,\mathrm{exneg}_P\,B\) means extreme positivity to \(\neg B\).
Conclusive relevance therefore includes both a conditional certainty
and a change from the prior, not merely certainty of the target.

For \(P(B)>0\), \((A\perp C\mid B)_P\) means
\(P(AC\mid B)=P(A\mid B)P(C\mid B)\).
When \(0<P(B)<1\), write \((A\perp C\mid\pm B)_P\) for the two
conditions given \(B\) and given \(\neg B\). No ratio conditional
on a zero-probability stratum is silently supplied.
\end{hbdefinition}

\begin{hbtheorem}{thm:p15-r43-polarity}{proved=annex:p15.r43.polarity}{Basic polarity laws and conclusive chains}
Relevance polarity is symmetric in its two arguments. Negating one
argument reverses the sign, while negating both preserves it. An event
of probability zero or one is irrelevant to every event, including itself.

Extreme positivity is transitive. If \(A\) logically entails \(B\),
\(P(A)>0\), and \(P(B)<1\), then \(A\,\mathrm{expos}_P\,B\).
The converse need not hold: an almost-sure implication can leave
probability-zero countercases in the underlying possibility space.

Ordinary positive relevance is not transitive. On three outcomes of
equal mass, let \(A\) be the first, \(C\) the second, and \(B=A\cup C\).
Then \(A\) is positive to \(B\), \(B\) positive to \(C\), but \(A\)
negative to \(C\).
\end{hbtheorem}

The three-outcome realization gives the source's counterexample its
explicit probability model. Merin illustrates the same form with the
2016-election propositions that Trump wins, Clinton wins, and a new
president is elected: the first two incompatible possibilities each
favor the third. These are the source's historical example propositions,
not a claim about a present election. His contrast with a Dow-rise
premise and the primality of seven is different: a target already assigned
one cannot receive an increment, even though its conditional probability
is also one. Merin relates extreme positivity to a colloquial conclusive
reason and indicative conditionals; that interpretive proposal is not
an identification with material implication in every model. His
\emph{P-entailment}, \(P(\neg A\vee B)=1\), expresses almost-sure
material implication; with positive antecedent mass and a prior target
below one it gives extreme positivity
\parencite[pp.~265--267]{Merin2021Relevance}.

\begin{hbtheorem}{thm:p15-r43-screening}{proved=annex:p15.r43.screening}{Conditional irrelevance, transitivity, and conjunction}
Let \(p=P(B)\in(0,1)\),
\(a_1=P(A\mid B),a_0=P(A\mid\neg B)\), and
\(c_1=P(C\mid B),c_0=P(C\mid\neg B)\).
Under \((A\perp C\mid\pm B)_P\),
\[
 \Delta_P(A,C)=p(1-p)(a_1-a_0)(c_1-c_0)
             =\frac{\Delta_P(A,B)\Delta_P(B,C)}{p(1-p)}.
\]
Thus the polarity between \(A,C\) is the product of the two link
polarities. If both links are positive, \(A,C\) are mutually positive
and \(AC\) is positive to \(B\). The positive-link transitivity
conclusion still holds if both within-stratum covariances are merely
nonnegative instead of zero.

If \(P(B\mid A)=1\), \(P(C\mid B)=1\), and \(0<P(B)<1\), with
the displayed ratios defined, then \((A\perp C\mid\pm B)_P\).
Consequently the conclusive chain is a special case of the screening
configuration, not a separate kind of correlation arithmetic.
\end{hbtheorem}

Conditional irrelevance is not unconditional irrelevance. In an editorial
numerical realization of the source's British/non-British drug-and-cure
example, assign the two groups equal mass. Within the first, let drug
administration and cure be independent with probabilities \(4/5,4/5\);
within the second, let them be independent with probabilities \(1/5,1/5\).
Then \(P(A)=P(C)=1/2\), \(P(AC)=17/50\), and
\(P(C\mid A)=17/25>1/2\), despite irrelevance within each group.
This is the source's Simpson-type contrast. It neither estimates a real
treatment effect nor proves causal ineffectiveness from conditional
independence alone.

Merin associates the screening configuration with Reichenbach's
common-cause interpretation and Pearl's suggestion that diagnostic
categories can be formed to make symptoms conditionally independent.
He suggests looking for counterexamples with a paradoxical conjunction
of symptoms. A failure to find one favors that thesis only under a
specified search/detection model; it does not itself establish the
categorization hypothesis
\parencite[p.~267 and p.~274, nn.~3--4]{Merin2021Relevance}.

\section{Predictions and the problem of old evidence}
\label{sec:p15-r43-old-evidence}

If \(H\subseteq D\) with \(0<P(H)\leq P(D)<1\), observing \(D\)
raises \(H\)'s probability to \(P(H)/P(D)\), while observing
\(\neg D\) reduces it to zero. Both evidence events have positive
probability. The prior uncertainty of \(D\) is indispensable. If
\(P(D)=1\), then \(P(H\mid D)=P(H)\): this already-known evidence
has no incremental relevance in that credal state. These consequences
are included in Theorem~\ref{thm:p15-r43-polarity}.

The source's old-evidence example concerns general relativity and the
long-known anomaly in Mercury's orbit. The reported increase of confidence
cannot be reconstructed as conditioning the same already-informed
state on that known fact again. Merin surveys hypothetical alternative
update histories, uncertainty about entailment/support, and a small
residual uncertainty paired with a suitable measure. He also distinguishes
a scientist's known data from a statistical model assigning distributions
to hypothetical repeatable observations. None of those approaches is
the same operation as setting \(P(D)=1\) and expecting an increment
from \(P_D\) without changing anything else
\parencite[pp.~267--268 and p.~274, n.~5]{Merin2021Relevance}.

\section{Directions, neutral values, and the principal measures}
\label{sec:p15-r43-measures}

Merin's notation \(R_T(A\Yright B)_P\) separates the \emph{source}
before the forked signpost from the \emph{target} after it. The separator
is not a logical implication. Assigning the names evidence and hypothesis
is optional: one can compare hypotheses as sources of possible data,
or compare propositions without that content classification. The notation
avoids confusing relevance with \(P(H\mid E)\) and exposes direction
changes that an unordered pair or an ambiguous comma can conceal.

His \emph{Transparency} requirement makes a measure depend on the joint
probabilities of \(EH,E\neg H,\neg EH,\neg E\neg H\), without
hidden propositions; differentiability is an interior regularity
condition. \emph{Quality Representation} assigns a neutral value
\(\iota_T\) and puts positive/negative relevance respectively above/below
it. Zero and one are the common additive and multiplicative neutrals.
Boundary definitions still have to be stated: differentiability does
not assign a value to \(0/0\)
\parencite[pp.~268--269 and p.~274, nn.~6, 8--9]{Merin2021Relevance}.

\begin{hbdefinition}{def:p15-r43-measures}{D, Q, L, their logarithms, KO, and Z}
For \(e=P(E)>0\), \(h=P(H)\), and \(q=P(H\mid E)\), define
\[
 R_D(E\Yright H)_P=q-h,
 \qquad R_Q(E\Yright H)_P=q/h\quad(h>0).
\]
For \(0<h<1\), put \(a=P(E\mid H)\), \(b=P(E\mid\neg H)\).
Then \(R_L(E\Yright H)_P=a/b\) when \(b>0\), and is \(+\infty\)
when \(b=0<a\). The source's completions are \(R_D=0\) if \(e=0\),
\(R_Q=1\) if \(e=0\) or \(h=0\), and \(R_L=1\) if the target is
certain or impossible, or both likelihoods vanish.

Define \(R_{QL}=\ln R_Q\) and \(R_{LL}=\ln R_L\), with
\(\ln0=-\infty\) and \(\ln(+\infty)=+\infty\). The expanded
log-probability differences apply only where their operands are defined;
the neutral completions give zero logs, not \(-\infty-(-\infty)\).

When \(0<e,h<1\), the source's bounded likelihood transform is
\[
 R_{KO}(E\Yright H)_P=\frac{a-b}{a+b},
\]
and its completion elsewhere is zero. Finally, on \(e>0,0<h<1\),
\[
 R_Z(E\Yright H)_P=
 \begin{cases}(q-h)/(1-h),&q>h,\\(q-h)/h,&q\leq h,\end{cases}
\]
with zero completion when evidence or target degeneracy prevents that
definition. These are mathematical conventions for neutral relevance,
not conditionals defined at impossible evidence.
\end{hbdefinition}

\begin{hbtheorem}{thm:p15-r43-measures}{proved=annex:p15.r43.measures}{Interpretations, symmetry, and negation}
On the defined nondegenerate domains,
\[
 \begin{aligned}
 q&=h+R_D=hR_Q,& O(H\mid E)&=O(H)R_L,\\
 \ln q&=\ln h+R_{QL},&
 \ln O(H\mid E)&=\ln O(H)+R_{LL},
 \end{aligned}
\]
where \(O(x)=x/(1-x)\) for a probability \(x\in(0,1)\), and
\(O(H)\) abbreviates \(O(P(H))\). Scalar endpoints, where used, are
\(O(0)=0\) and \(O(1)=+\infty\). Log identities use positive finite operands
or their explicitly valid one-sided limits. Q and QL are amount-symmetric.
KO is \(\tanh(\tfrac12\ln R_L)\) and is ordinally equivalent to L
and LL on their applicable domain. D, L, KO, and Z preserve the sign
of qualitative relevance relative to their stated neutral.

For fixed \(h\in(0,1)\), Z is continuous and strictly increasing
in the posterior \(q\), but is differentiable at \(q=h\) only when
\(h=1/2\). Joint-table transparency must therefore not be conflated
with global differentiability of every listed measure.

Negating the target negates D and Z and inverts L; it negates LL.
Negating the source need not have these same numerical effects. If
\(0<P(A),P(B)<1\), the relevance is nonzero, and
\(\delta=\Delta_P(A,B)\), then
\[
 |R_D(A\Yright B)_P|=|\delta|/P(A),
 \quad |R_D(B\Yright A)_P|=|\delta|/P(B).
\]
Thus the less probable source has the larger absolute D relevance.
L has the corresponding direction comparison away from its neutral,
apart from the tied zero values of extreme negative relevance.
\end{hbtheorem}

D has range \((-1,1)\), Q has \([0,\infty)\), L has
\([0,+\infty]\), QL has \([-\infty,\infty)\), LL has
\([-\infty,+\infty]\), and KO and Z have \([-1,1]\), with the
domain qualifications just given. L and Z reach their extremes exactly
at conclusive positive or negative relevance on a nondegenerate target.
D instead assigns the same increase \(0.2\) to the source's shifts
\(0.2\to0.4\), \(0.5\to0.7\), and \(0.8\to1\). A shift
\(0.999\to1\) has infinite L, whereas \(0.001\to0.999\) has the
finite L value \(998001\). Conclusiveness and amount of probability
change therefore need not induce the same ordering.

The source's divisibility example illustrates direction: knowing a
natural number at least fifteen is divisible by fifteen entails
divisibility by three, not conversely. If the probability of the first
event is strictly between zero and that of the second, D distinguishes
the two directions, while Q treats them alike. The distribution over
numbers must be supplied; a uniform probability on all such natural
numbers is not silently assumed
\parencite[pp.~269--272 and p.~274, nn.~7, 10, 17]{Merin2021Relevance}.

\begin{hbtheorem}{thm:p15-r43-sufficiency}{proved=annex:p15.r43.sufficiency}{Likelihood ratio as a minimal sufficient statistic in a finite binary experiment}
Let two simple hypotheses give masses \(f_1(x),f_0(x)\) on a finite
data space, and retain only \(x\) with \(f_1(x)+f_0(x)>0\).
The extended likelihood ratio \(\Lambda(x)=f_1(x)/f_0(x)\), with
\(\Lambda=+\infty\) when only \(f_1\) is positive, is sufficient:
within each value class the conditional data law does not depend on
which hypothesis generated it whenever that conditional law is defined.
Every sufficient statistic refines these classes, so \(\Lambda\)
is a function of any such statistic. Its logarithm has the same classes.
\end{hbtheorem}

This fixes the scope of Merin's sufficient-statistic reason for preferring
L or LL: two specified alternative data laws. A binary-hypothesis
likelihood ratio does not establish sufficiency for an arbitrary
multi-parameter family. In Bayesian Model Comparison,
Definition~\ref{def:p15-model-comparison-experiment}, analogous factors
use prior-predictive laws, after integrating parameter priors. They
remain distinct from maximized likelihoods and posterior model weights
\parencite[p.~270 and p.~274]{Merin2021Relevance}.

\section{Axioms characterize different ordinal families}
\label{sec:p15-r43-axioms}

Ordinal equivalence means preserving all order comparisons and ties.
Fixed strictly increasing transforms therefore preserve rankings, though
they need not preserve addition or multiplication. Merin distinguishes
four axioms, using \(\bowtie\) for any of \(<,=,>\):

\begin{enumerate}
\item POP-D: for a fixed target \(H\),
\(R(E_1\Yright H)_P\bowtie R(E_2\Yright H)_P\) exactly when
\(P(H\mid E_1)\bowtie P(H\mid E_2)\).
\item LIK-D: for fixed source \(E\), comparisons of
\(R(E\Yright H_i)_P\) follow comparisons of \(P(E\mid H_i)\).
The targets need not be incompatible or exhaustive.
\item MCIE: if \((E_1\perp E_2\mid\pm H)_P\), updating on \(E_1\)
leaves the relevance amount of \(E_2\) to \(H\) unchanged.
The prior and updated target must remain nondegenerate for this
axiomatic comparison; neutrality completions at dogmatic targets do
not make it an unrestricted assertion.
\item MDHD: for \(P(H_1H_2)=0\),
\(R(E\Yright H_1\vee H_2)_P\bowtie R(E\Yright H_1)_P\) exactly
when \(P(H_2\mid E)\bowtie P(H_2)\), on the comparison domain.
\end{enumerate}

\begin{hbtheorem}{thm:p15-r43-characterization}{proved=annex:p15.r43.characterization}{The Big Three's ordinal characterization on defined comparison domains}
Suppose one formal function of the evidence--target joint table is used
uniformly over all finite probability models. Restrict its comparisons
to evidence and target probabilities strictly between zero and one;
updated contexts in MCIE have this same target domain. No smoothness
is needed. Then:
\[
 \begin{array}{c|c}
 \text{additional axioms}&\text{form of the measure}\\\hline
 \text{POP-D and MDHD}&f_D(q-h)\\
 \text{POP-D and LIK-D}&f_Q(q/h)\\
 \text{POP-D and MCIE}&f_L\bigl(O(q)/O(h)\bigr)
 \end{array}
\]
Each \(f\) is strictly increasing on the attained argument domain,
and each displayed family conversely satisfies its axioms.
For L, zero and infinite arguments have ordered endpoint values in
the chosen output range; log arithmetic is not used to define an
indeterminate expression.
These axioms do not impose arbitrary neutral completions outside their
comparison domain.
\end{hbtheorem}

The full construction behind the characterization is supplied in the
annex, using the actual proof dependencies
\parencite[pp.~192--201]{CrupiChaterTentori2013Axioms}
\parencite[pp.~654--656, 662--665]{CrupiTentori2016Confirmation}.
It compares realizable joint models, not arbitrary inconsistent triples.
The resulting equivalence classes do not single out a unit, neutral
encoding, or numerical transform.

Merin judges POP-D particularly attractive for first-resort measures.
Accepting it together with amount symmetry also gives LIK-D, by applying
POP-D to the reversed comparisons. Users who want asymmetric valuation
of conclusive reasons therefore cannot demand that entire combination.
MCIE offers useful composition; MDHD has a different cost. If E rules
out H conclusively, but G is disjoint from H and irrelevant to E, then
D gives H and \(H\vee G\) the same decrement, despite the latter not
being ruled out. For example, in a fair fifty-two-card model, a face-card
observation excludes the seven of spades, whose prior is \(1/52\),
while leaving the probability of a red card at \(1/2\). For H the
seven of spades and G red, both H and \(H\vee G\) have D value
\(-1/52\); Q values are respectively zero and \(26/27\).
This illustrates the source's MDHD objection, not a contradiction
within D's definition
\parencite[pp.~270--271 and p.~274, nn.~11--14]{Merin2021Relevance}.

\section{Which rankings are robust, and which are not?}
\label{sec:p15-r43-rankings}

\begin{hbtheorem}{thm:p15-r43-rankings}{proved=annex:p15.r43.rankings}{Fixed-target agreement, target disagreement, and direction}
For one fixed target \(H\), every measure satisfying POP-D gives
the same ranking and ties among defined evidence sources \(E_i\).
This includes D, Q, L and their increasing transforms, and Z.

For a fixed evidence event, D, Q, and L can disagree pairwise about
different targets. The following two realizable comparison models use
\(P(E)=1/5\):
\[
 \begin{array}{c|cc|ccc}
 &P(H)&P(H\mid E)&D&Q&L\\\hline
 \text{I: }H_1&3/5&9/10&3/10&3/2&6\\
 \text{I: }H_2&1/10&3/10&1/5&3&27/7\\
 \text{II: }H_1&1/5&1/2&3/10&5/2&4\\
 \text{II: }H_2&4/5&19/20&3/20&19/16&19/4
 \end{array}
\]
Model I separates D from Q and Q from L; model II separates D from L.

Swapping source and target can also change a ranking. With
\(P(H)=1/2\), let \(P(E_1)=1/20,P(H\mid E_1)=4/5\), and
\(P(E_2)=1/2,P(H\mid E_2)=7/10\). Forward D values are
\(3/10,1/5\), while reversed D values are \(3/100,1/5\), giving
opposite rankings. Q is symmetric and preserves its ranking in either
direction. It can therefore implement the fixed-target ranking without
directional ambiguity; this does not make every measure direction-invariant.
\end{hbtheorem}

For Merin's clouds-versus-forecast comparison concerning one rain target,
POP-D supplies useful measure robustness. Comparing rain and snow as
different targets, or simply reversing argument positions in an
asymmetric measure, is a different task. None is answered by the
qualitative sign alone
\parencite[pp.~268, 271--272]{Merin2021Relevance}.

The raven paradox provides a source example. A black raven can favor
the universal raven hypothesis H; a nonblack nonraven can favor its
logically equivalent contrapositive G. But logical equivalence alone
does not determine the sizes of those increments. Merin invokes
Hosiasson-Lindenbaum's independence assumptions about the separately
observed property raven and the property nonblack, as well as their
very different background frequencies.

In an explicit model, let \(P(H)=1/2\). Conditional on H, give raven
probability \(1/100\), nonblack probability \(4/5\), and their
intersection zero. Conditional on \(\neg H\), retain the two
marginals but give their intersection \(1/200\). These are valid
joint tables because the remaining cells have nonnegative mass.
For A a black raven and B a nonblack nonraven,
\[
 P(H\mid A)=2/3,
 \qquad P(H\mid B)=160/319,
 \qquad L(A\Yright H)=2,
 \quad L(B\Yright H)=160/159.
\]
Both favor H, but the second does so much less. Every POP-D measure
agrees on this ordering. A white shoe entails the specimen description
nonblack nonraven, not a universal claim about all nonblack things.
Its own relevance requires a further observation model, such as
conditional independence of the shoe subtype and H within B. Entailment
to B alone does not supply that refinement assumption. These calculations
and that distinction are proved in Annex~\ref{annex:p15.r43.rankings}
\parencite[pp.~271--272 and p.~274, nn.~15--16]{Merin2021Relevance}.

\section{Conjunction, disjunction, and their different weights}
\label{sec:p15-r43-composition}

\begin{hbtheorem}{thm:p15-r43-composition}{proved=annex:p15.r43.composition}{Compound-source relevance under its actual assumptions}
For \(0<P(H)<1\), conditionally independent evidence given both
H and \(\neg H\) has
\[
 L(E_1E_2\Yright H)=L(E_1\Yright H)L(E_2\Yright H),
 \quad LL(E_1E_2\Yright H)=LL(E_1\Yright H)+LL(E_2\Yright H)
\]
when the likelihoods are positive and finite. An arbitrary increasing
transform of L need not preserve these operations.

If \(P(AB)=0\), \(P(A)+P(B)>0\), and the target is nondegenerate,
D and Q relevance of \(A\vee B\) is the weighted mean of their
relevances with weight \(\alpha=P(A)/(P(A)+P(B))\).
For finite L values with positive denominator likelihoods, its weight
is instead \(\alpha=P(A\mid\neg H)/(P(A\mid\neg H)+P(B\mid\neg H))\).
Thus the mixture lies between constituent values, but different
measures need not use the same weights. Increasing transforms preserve
betweenness, not those fixed arithmetic-mean formulas.

For Carnap's covariance measure \(R_C(E\Yright H)=\Delta_P(E,H)\),
disjoint source disjunction is additive, and source conjunction is
additive when \(P(A\vee B)=1\). This measure does not obey POP-D.

Under conditional independence given \(\pm H\), if both A and B have
finite L strictly greater than one, then
\[
 1<L(A\vee B\Yright H)<\max\{L(A\Yright H),L(B\Yright H)\}
                   <L(AB\Yright H).
\]
For \(A\mathbin{\dot\vee}B=(A\vee B)\wedge\neg(AB)\),
\(L(A\mathbin{\dot\vee}B\Yright H)<L(A\vee B\Yright H)\),
but the exclusive disjunction can be positive, neutral, or negative.
The same fixed-target order comparisons hold for every POP-D measure.
\end{hbtheorem}

Merin's note~16 uses the last comparison to explain why adding a
defeasible not-both implicature to inclusive \emph{or} can reduce
evidential relevance. The first disjunction rule has a different
premise, disjointness, and must not be substituted for conditional
independence. In the conditionally independent configuration, inclusive
or and and generally overlap. Merin's quasi-utility interpretation
describes relevance as a resource toward establishing a target in a
credal/rhetorical context. The shared mixture form does not identify
relevance with welfare, a person's utilities, or a decision rule
\parencite[pp.~271--274 and p.~274, n.~16]{Merin2021Relevance}.

\section{Psychological interpretation and reference use}
\label{sec:p15-r43-use}

Merin suggests that some apparent base-rate-neglect and conjunction
responses may be relevance judgments offered in place of requested
probability judgments. The arithmetic distinction is exact: replacing
prior odds by one makes posterior odds equal L, although that equality
does not hold with other prior odds. In his Linda example, participants
are asked to rank bank teller, feminist, and feminist bank teller after
a personality description. Merin reprises a fictional thirty-one-year-old
philosophy graduate whose description stresses intellectual ability,
forceful speech, unmarried status, and equal-opportunity activism.
He treats that description as favoring feminist activity and disfavoring
bank-teller work. The probability law
\(P(BF\mid E)\leq P(B\mid E)\) cannot be reversed. A relevance
comparison can reverse without violating that law, because the prior
and posterior for different targets both matter. For example, in an
editorial model with \(P(E)=1/10\), give the four cells
\(BF,B\neg F,(\neg B)F,\neg B\neg F\) prior masses
\((.02,.18,.48,.32)\) and conditional masses given E
\((.10,.05,.70,.15)\). The conjunction posterior is .10, below the
bank-teller posterior .15, but its D relevance is .08, above the
bank-teller relevance -.05. Feminist activity rises from .50 to .80.
These are coherent stipulated values, not participant measurements;
the realization is checked in Annex~\ref{annex:p15.r43.rankings}.

This possible reinterpretation is a psychological hypothesis, not
proof of what a respondent intended. The chapter points to Tentori,
Crupi and Russo, and to Cevolani--Crupi, but supplies no participant
counts or effect estimates. The actual original conjunction tasks and
language/response comparisons remain in
Chapters~\ref{sec:p19-conjunction-judgments} and
\ref{sec:p19-conjunction-language}. No new population-frequency claim
or empirical completion credit is inferred here.

Merin distinguishes Sperber--Wilson's linguistic relevance theory,
concerned with cognitive effects and processing costs, from probabilistic
evidential relevance. He identifies the Wason-task application and
Skovgaard-Olsen--Singmann--Klauer's conditional/reason-relation research,
including \(\delta p=P(H\mid E)-P(H\mid\neg E)\).
When \(0<P(E)<1\), this latter quantity equals
\(R_D(E\Yright H)/(1-P(E))\), by total probability, so it is not
the same numerical measure as D. The chapter provides research leads,
not an empirical record sufficient to establish universal covert
relevance processing or a uniquely used human measure
\parencite[pp.~273--274 and p.~274, nn.~18--20]{Merin2021Relevance}.

\begin{hbmethod}{def:p15-r43-selection}{Selecting a relevance comparison}
State the probability model, background information, evidence source,
target, and intended conditioning regime. Separate the sign of the
increment, its numerical size, final posterior confidence, and the
decision criterion. Check support before evaluating ratios and logs.

For several sources and one fixed target, say whether POP-D is required;
if so, ordinal rankings are robust across the admitted measures.
For several targets, declare the measure and the role of priors.
Do not reverse asymmetric arguments without recalculating.

For compound evidence, establish the relevant premise: conditional
independence under both alternatives, disjoint source events, or another
joint observation model. Keep the weights and operation appropriate
to the chosen measure. Duplicate reports supply no new information
solely by receiving new names.

For a claim about people, specify the actual requested quantity,
response interpretation, task, population, design, and uncertainty.
A mathematically possible relevance ranking does not prove that
participants performed it. For a claimed logical consequence or causal
effect, provide the additional semantic or intervention premises.
\end{hbmethod}

The posterior-model-odds result in
Theorem~\ref{thm:p15-model-posterior-odds} and the P18 detection example
in Theorem~\ref{thm:p18-r55-detection} give exact neighboring uses:
a factor or increment is not a posterior, and high confidence can
reflect the prior alone. Argument schemes and dialogue obligations
have different reconstruction and procedural purposes. Merin recommends
LL as a practical default for quantitative evidential work because of
its sufficient-statistic and compositional properties, Q/QL when
direction is unspecified, and explicit pluralism where other
desiderata matter. These are defended preferences about use, not
empirical consensus or consequences of a symbol's name
\parencite[p.~274]{Merin2021Relevance}.

\index{evidential relevance}\index{confirmation!incremental}
\index{likelihood ratio!relevance}\index{conditional independence!relevance}
\index{evidence!compound}\index{raven paradox}\index{old evidence}
